\chapter{Запис і читання}
% full version: chapters/09_acet.tex

Пам'ять у мозку має незручну особливість: спогади зберігають ті самі зв'язки, якими йде робота. Окремого жорсткого диска немає. Щоразу, коли мережа згадує, вона користується своїми зв'язками, --- і щоразу, коли вчиться, вона їх змінює. Як не зіпсувати старе, записуючи нове? У комп'ютері для цього є окремий сигнал --- «дозволити запис». У мозку його роль, схоже, грає радіостанція \nt{ACET} --- ацетилхолін.

\section*{Два режими}

Уявіть мережу, яка зберігає кілька картинок: покажіть їй половину --- вона домалює решту. Це читання. Тепер покажіть їй нову картинку, схожу на стару. Якщо мережа за звичкою домальовуватиме, вона побачить не нове, а старе --- і вивчить помилку. Щоб записати нове, мережа має на час перестати довіряти своїй пам'яті й дивитися на світ.

\pencilfig{9}{\pencilart{9}}{Один провід перемикає мережу: записувати нове чи згадувати старе.}

Ацетилхолін робить саме це, і всі три його ефекти виміряно: він послаблює внутрішні зв'язки мережі, через які вона «згадує», підсилює сигнали, що надходять від органів чуттів, і полегшує зміну зв'язків. Багато ацетилхоліну --- режим запису: мережа слухає світ і вчиться. Мало --- режим читання: мережа спирається на пам'ять і домальовує.

У нашій моделі немає такого рівня, за якого добре йшло б і те, і інше. Для запису внутрішні зв'язки треба приглушити, а для читання вони потрібні на повну силу. Тому перемикач необхідний.

\section*{Скільки довіряти очам}

Інженери давно розв'язали схожу задачу. Коли корабель іде курсом, у штурмана є розрахунок, де корабель має бути, і є нові, але неточні спостереження. Скільки довіряти кожному? Оптимальний фільтр відповідає: довіряй спостереженням тим більше, чим менше впевнений у розрахунку. І в цьому ж рецепті є друге: що менша впевненість, то швидше треба перевчатися. Одна ручка керує водночас і вагою входу, і швидкістю навчання. Це саме те, що робить ацетилхолін. Звідси гіпотеза: він повідомляє мережі, наскільки непевна її власна модель світу.

\pencilfig{124}{%
% optimal blending: the less sure the model, the more weight on the observation (and the faster it relearns)
\foreach \dx/\wm/\we/\ttop/\bbot in {0/2.4pt/0.6pt/{модель упевнена}/{ацетилхоліну мало}, 4.3/0.6pt/2.4pt/{модель не впевнена}/{ацетилхоліну багато}}{
  \begin{scope}[xshift=\dx cm]
  \draw[penlite=pcViolet, wash=pcViolet, rounded corners=3pt] (0,0.95) rectangle (1.3,1.45); \node[lbl, text=black] at (0.65,1.2) {пам'ять};
  \draw[penlite=pcBlue, wash=pcBlue, rounded corners=3pt] (0,-0.25) rectangle (1.3,0.25); \node[lbl, text=black] at (0.65,0) {очі};
  \draw[dotpen=pcAmber, wash=pcAmber] (2.95,0.6) circle (0.55); \node[lbl, text=black] at (2.95,0.6) {оцінка};
  \draw[pcViolet, line width=\wm, line cap=round, -{Stealth[length=5pt]}] (1.35,1.2) -- (2.4,0.8);
  \draw[pcBlue, line width=\we, line cap=round, -{Stealth[length=5pt]}] (1.35,0) -- (2.4,0.4);
  \node[font=\small\itshape, text=pcGray, anchor=south] at (1.55,1.6) {\ttop};
  \node[lbl, anchor=north] at (1.55,-0.4) {\bbot};
  \end{scope}}
}{Скільки довіряти очам. Якщо мережа впевнена у своїй моделі світу, оцінка майже вся з пам'яті. Якщо не впевнена --- зі спостережень, і тоді ж мережа швидше перевчається. За гіпотезою, цю ручку крутить ацетилхолін.}

\section*{Поділ праці}

Згадайте попередній розділ. Коли світ раптово змінюється, норадреналін, за гіпотезою, помічає несподіванку. А ацетилхолін тримає мережу в режимі запису, доки нову обстановку не вивчено. Один повідомляє «щось змінилося», інший --- «пиши».

\section*{Нічний перезапис}

Під час глибокого сну рівень ацетилхоліну низький. Мережа відключена від світу й перебуває в режимі читання. Саме в цей час у ній «програються» денні події --- це виміряно. Що при цьому вивчене за день переписується в довготривалу пам'ять, --- правдоподібна гіпотеза, до якої ми повернемося в розділі про сон.

\fullversion{9}
